% Reviewed source SHA256 (UTF-8/LF): 222de12c79ca779329500977da68fb58ca3f0a10a0ae11058550db303ac8fe8b
\documentclass[11pt]{article}
\usepackage[T1]{fontenc}
\usepackage{lmodern}
\usepackage[a4paper,margin=27mm]{geometry}
\usepackage{amsmath,amssymb,amsthm,mathtools}
\usepackage{microtype,booktabs,array}
\usepackage[hidelinks]{hyperref}
\usepackage{enumitem}
\setlist{itemsep=3pt,topsep=5pt}
\newtheorem{theorem}{Theorem}[section]
\newtheorem{lemma}[theorem]{Lemma}
\newtheorem{corollary}[theorem]{Corollary}
\newtheorem{proposition}[theorem]{Proposition}
\theoremstyle{definition}
\newtheorem{definition}[theorem]{Definition}
\newtheorem{remark}[theorem]{Remark}
\newcommand{\R}{\mathbb R}
\newcommand{\E}{\mathbb E}
\newcommand{\Prob}{\mathbb P}
\newcommand{\tr}{\operatorname{tr}}
\newcommand{\rank}{\operatorname{rank}}
\newcommand{\diag}{\operatorname{diag}}
\newcommand{\range}{\operatorname{range}}
\newcommand{\OPT}{\operatorname{OPT}}
\newcommand{\norm}[1]{\left\lVert#1\right\rVert}
\newcommand{\ip}[2]{\left\langle#1,#2\right\rangle}
\newcommand{\psd}{\succeq}
\newcommand{\pd}{\succ}
\newcommand{\eps}{\varepsilon}
\DeclareMathOperator*{\argmin}{arg\,min}
\allowdisplaybreaks[1]
\hypersetup{pdfcreator={LaTeX},pdfauthor={Matthew J. Colbrook}}

\title{Counterexamples to two auxiliary Gamma-density conjectures}
\author{Matthew J. Colbrook\\\small Department of Applied Mathematics and Theoretical Physics\\\small University of Cambridge, Cambridge, United Kingdom\\\small\href{mailto:m.colbrook@damtp.cam.ac.uk}{m.colbrook@damtp.cam.ac.uk}}
\date{11 September 2026}
\usepackage{xurl}
\setlength{\emergencystretch}{3em}
\begin{document}
\maketitle

\noindent\textbf{Verified auxiliary counterexamples.} Propositions 2.1 and 3.1 refute the auxiliary upper-mode and upper-inflection assertions. The catalog trace-tail questions RA-12 and RA-13 remain open.
Independent agent verification is not external human peer review or formal certification. Original draft remarks about review or source access reflect the input date; the dated \href{https://github.com/MColbrook/OpenProblemsInNLA/blob/codex/colbrook-transfer-resolutions/references/colbrook-transfer-2026-09-11/verification/reviews/gamma-auxiliary-review.md}{independent review} records the subsequent checks.
\par\medskip
% BEGIN REVIEWED BODY
\begin{abstract}
We refute Conjectures~1 and~2 of Hallman's \emph{Extremal bounds for Gaussian trace estimation}, arXiv:2411.15454v1. The counterexamples use four independent Gamma variables with common shape $1/2$, two independent exponential variables, and an elementary convolution density. Its first derivative is positive beyond the conjectured mode bound, and its second derivative changes sign beyond the conjectured inflection-point bound. These are counterexamples to the auxiliary density-location conjectures. They do not by themselves refute the distinct extremal trace-tail conjectures numbered~3 and~4 in that paper.
\end{abstract}

\section{The distribution}
Use the shape--rate convention: $\Gamma(\alpha,\beta)$ has mean $\alpha/\beta$ and variance $\alpha/\beta^2$. Let
\[
 G_1,G_2,G_3,G_4\ \text{be independent }\Gamma(1/2,1),
 \qquad Z_1,Z_2\ \text{be independent }\Gamma(1,1),
\]
with all six variables mutually independent. Set
\begin{equation}\label{eq:QY}
 Q=G_1+G_2+\tfrac12(G_3+G_4),\qquad Y=Q+Z_1+Z_2.
\end{equation}
Then
\[
 \E Q=\frac32,\qquad \operatorname{Var}(Q)=\frac54,
 \qquad Y\ \stackrel{d}=\ \Gamma(3,1)+\Gamma(1,2),
\]
where the two variables on the right are independent. Direct convolution gives, for $x>0$,
\begin{align}\label{eq:density}
 p(x)&=\int_0^x\frac{(x-y)^2}{2}e^{-(x-y)}\,2e^{-2y}\,dy\\
 &=e^{-x}(x^2-2x+2)-2e^{-2x}.\nonumber
\end{align}
Consequently,
\begin{align}
 p'(x)&=-e^{-x}(x-2)^2+4e^{-2x},\label{eq:first}\\
 p''(x)&=e^{-x}(x-2)(x-4)-8e^{-2x}.\label{eq:second}
\end{align}

\section{The mode conjecture}
Hallman \cite[Theorem~2 and Conjecture~1]{Hallman2024} considers nonnegative sums $Q_\lambda=\sum_i\lambda_iG_i$ with common Gamma shape $\alpha$ and rate $\beta$, mean one, and scale $\max_i\lambda_i/\beta\le1/\mu$. Two independent exponentials of rate $\beta$ are added with distinct coefficient indices $j\ne k$. The conjectured upper bound for every resulting mode is $1+1/\mu$.

For \eqref{eq:QY}, normalize by the mean of $Q$:
\[
 Q'=\tfrac23Q,
 \qquad \lambda=(2/3,2/3,1/3,1/3),
 \qquad \alpha=1/2,\quad\beta=1,\quad\mu=3/2.
\]
The mean of $Q'$ is one, and its scale is $2/3=1/\mu$, so it belongs to the stipulated family. Taking $j=1$ and $k=2$ gives the legally indexed augmented variable
\[
 Y'=Q'+\tfrac23 Z_1+\tfrac23 Z_2=\tfrac23Y.
\]
The proposed upper-mode bound $1+1/\mu=5/3$ for $Y'$ is equivalent to an upper-mode bound $5/2$ for $Y$.

\begin{proposition}
The unique positive mode of $Y$ exceeds $5/2$. Hence the upper bound in Hallman's Conjecture~1 is false.
\end{proposition}
\begin{proof}
At the proposed bound,
\[
 p'(5/2)=e^{-5/2}\bigl(4e^{-5/2}-1/4\bigr)>0.
\]
For an elementary exact sign certificate, $e<3$ implies $e^{5/2}<9\sqrt3<16$, with the last inequality equivalent to $243<256$.

We also verify that this is the increasing side of the unique mode, rather than relying only on one derivative value. For $0<x\le2$, strict convexity of the exponential gives $2e^{-x/2}>2-x$, so \eqref{eq:first} is positive. For $x>2$, its unique zero is determined by
\[
 x-2=2e^{-x/2}.
\]
The left side is strictly increasing and the right side strictly decreasing. The derivative is positive before this zero and negative afterward, so it is the unique mode. Since $p'(5/2)>0$, the mode exceeds $5/2$.
\end{proof}

Equivalently, the mode is $2+2W(e^{-1})$, where $W$ is the real inverse of $z\mapsto ze^z$ on $[0,\infty)$. Numerically,
\[
 \operatorname{mode}(Y)\approx2.55692908552,
 \qquad\operatorname{mode}(Y')\approx1.70461939035>5/3.
\]
Only the upper half of Conjecture~1 is refuted here; no counterexample to its separately stated lower bound is claimed.

\section{The inflection-point conjecture}
Hallman \cite[Theorem~4 and Conjecture~2]{Hallman2024} fixes a scale bound $\lambda$ and variance $\phi^2$, and considers all inflection points of
\[
 Q_\lambda+\lambda_jZ_1+\lambda_kZ_2-\E Q_\lambda,
 \qquad j\ne k.
\]
The proposed upper bound is $\lambda+\sqrt{\phi^2+\lambda^2}$. Use the unnormalized $Q$ in \eqref{eq:QY}, with coefficient vector $(1,1,1/2,1/2)$, $\alpha=1/2$, $\beta=1$, $\lambda=1$, and $\phi=\sqrt5/2$. All family constraints hold, including $\lambda\le\phi/\sqrt\alpha$. With $j=1$ and $k=2$, the centered augmented variable is $Y-3/2$.

\begin{proposition}
The density of $Y-3/2$ has an inflection point strictly larger than $5/2=1+\sqrt{5/4+1}$. Therefore Hallman's Conjecture~2 is false.
\end{proposition}
\begin{proof}
The conjectured endpoint $5/2$ for $Y-3/2$ corresponds to $x=4$ for $Y$. But
\[
 p''(4)=-8e^{-8}<0,
 \qquad p''(5)=e^{-10}(3e^5-8)>0.
\]
The second sign follows already from $e^5>1+5=6$. Thus $p''$ changes sign somewhere in $(4,5)$. More explicitly, for $x>4$ its zeros satisfy
\[
 (x-2)(x-4)e^x=8.
\]
The left side is strictly increasing from zero to infinity, so there is exactly one such zero, and it is a genuine inflection point. Subtracting the centering constant $3/2$ puts it strictly above $5/2$.
\end{proof}
Numerically, the upper inflection point of $Y$ is about $4.06635720605$, and that of $Y-3/2$ is about $2.56635720605$.

\section{What these examples do and do not establish}
All weights are positive and the common base shape is $1/2$, so the examples occur within the Gamma representation of one-sample Gaussian quadratic forms. The two exponential increments use distinct indices, even though the corresponding coefficients are equal. Both features are allowed in the definitions being refuted.

The auxiliary quantities in Theorems~2 and~4 are suprema over broad families of augmented densities. The later Conjectures~3 and~4 instead concern particular final extremal trace-tail comparisons. Failure of an auxiliary bound does not logically imply failure of either final comparison. This note asserts no resolution of those two later conjectures, nor of a repository item that asks only for those final comparisons.

The exact convolution, derivatives, normalization and sign certificates are reproduced in \texttt{verify\_counterexamples.py}. Decimal roots are included for orientation only. This note is a candidate submission requiring independent review; it does not claim verified priority over all subsequent work.

\begin{thebibliography}{9}
\bibitem{Hallman2024}
E.~Hallman,
\emph{Extremal bounds for Gaussian trace estimation},
\href{https://arxiv.org/abs/2411.15454}{arXiv:2411.15454}, version~1, November~23, 2024.
The relevant statements are Theorem~2 and Conjecture~1 on page~9, and Theorem~4 and Conjecture~2 on page~10 of the PDF.
\end{thebibliography}
\end{document}

% END REVIEWED BODY
